%% ============================================================
%%  SECCIÓN XII — Observaciones
%% ============================================================

\section{Observaciones}

Durante el desarrollo del experimento se identificaron las siguientes
fuentes de error e incidencias que pudieron afectar la precisión de
los resultados:

\begin{enumerate}

    \item \textbf{Error de paralaje en la lectura de la regla:}
    Al medir la longitud del resorte y la liga suspendidos verticalmente,
    la posición del observador respecto a la escala graduada introduce
    un error sistemático difícil de eliminar por completo, a pesar de
    las precauciones adoptadas.

    \item \textbf{Oscilaciones del sistema:}
    Al colocar o retirar las masas, el resorte y la liga experimentaron
    oscilaciones transitorias. Aunque se esperó al equilibrio antes de
    registrar cada medición, no puede descartarse que algunos valores
    hayan sido tomados con una deformación ligeramente distinta a la
    de equilibrio estático.

    \item \textbf{Deformación remanente de la liga:}
    Se observó que, al retirar todas las masas, la liga no recuperó
    exactamente su longitud natural original, evidenciando una
    deformación plástica parcial acumulada a lo largo del experimento.
    Este efecto contribuye a los valores de histéresis registrados.

    \item \textbf{Precisión limitada de la regla métrica:}
    Con una resolución de \qty{1}{mm} e incertidumbre de
    \qty{\pm 0{,}5}{mm}, la regla constituye la principal fuente de
    error aleatorio en las mediciones de longitud, especialmente para
    la liga, cuyas deformaciones iniciales son pequeñas en comparación con su longitud natural.

    \item \textbf{Instrumento compartido:}
    La balanza fue utilizada por varios grupos simultáneamente, lo que
    limitó la posibilidad de verificar la masa de cada configuración
    con la frecuencia deseada, dependiendo de los valores nominales en algunos casos.

    \item \textbf{Alineación vertical:}
    La liga, por su menor rigidez lateral respecto al resorte, mostró
    mayor tendencia a desviarse de la vertical bajo carga, lo que pudo
    introducir una componente horizontal no considerada en el modelo de tracción pura.

\end{enumerate}

%% ============================================================
%%  SECCIÓN XIII — Sugerencias
%% ============================================================

\section{Sugerencias}

Con el fin de mejorar la precisión y reproducibilidad del experimento
en sesiones futuras, se proponen las siguientes recomendaciones:

\begin{enumerate}

    \item \textbf{Emplear un sensor de posición o extensómetro:}
    Sustituir la regla métrica por un sensor digital de desplazamiento
    reduciría significativamente el error de paralaje y permitiría
    registrar la deformación en tiempo real durante la carga y descarga.

    \item \textbf{Utilizar una cámara fija como referencia:}
    Montar una cámara alineada horizontalmente con la escala graduada
    permitiría realizar las lecturas de longitud \textit{a posteriori}, eliminando
    el error de paralaje del operador mediante software de análisis de imagen.

    \item \textbf{Permitir mayor tiempo de estabilización:}
    Especialmente para la liga, establecer un tiempo mínimo de espera
    fijo (por ejemplo, \qty{30}{s}) antes de registrar cada medición
    garantizaría lecturas en equilibrio estático real, mitigando efectos viscoelásticos transitorios.

    \item \textbf{Limitar el rango de carga de la liga:}
    Para evitar la deformación plástica parcial observada, conviene
    determinar previamente el límite elástico del material elastomérico y no superarlo
    durante el experimento.

    \item \textbf{Disponer de una balanza por grupo:}
    Contar con una balanza exclusiva por grupo eliminaría los tiempos
    de espera y permitiría verificar las masas exactas antes de cada etapa
    de carga.

    \item \textbf{Repetir el ciclo de carga y descarga:}
    Realizar al menos tres ciclos completos de carga–descarga y promediar
    los resultados permitiría estimar la incertidumbre estadística de
    los datos y cuantificar con mayor rigor el efecto de histéresis acumulada.

\end{enumerate}

%% ============================================================
%%  SECCIÓN XIV — Conclusiones
%% ============================================================

\section{Conclusiones}

A partir del análisis experimental y del tratamiento de datos realizados empleando el modelo de esfuerzo real, se establecen las siguientes conclusiones:

\begin{enumerate}

    \item \textbf{Verificación de la Ley de Hooke:}
    El resorte metálico exhibió una relación lineal entre la fuerza
    aplicada $F$ y la deformación longitudinal $\Delta\ell$ en el
    rango ensayado ($R^2 = 0{,}9966$), confirmando su comportamiento dentro de la región elástica. La constante recuperadora obtenida fue $k = \qty{80{,}89 \pm 2{,}38}{N/m}$.

    \item \textbf{Determinación del módulo de Young efectivo:}
    El módulo de Young del resorte, obtenido a partir del diagrama de esfuerzo real ($\sigma$) vs. deformación ($\varepsilon$), resultó ser $Y_r = \qty{90{,}23 \pm 2{,}65}{kPa}$. Este valor es órdenes de magnitud inferior al del acero masivo, lo cual es físicamente consistente dado que la geometría helicoidal disipa la carga mediante torsión y flexión del alambre.

    \item \textbf{Comportamiento no lineal del elastómero:}
    La liga elástica no siguió la Ley de Hooke; su diagrama $\sigma$--$\varepsilon$ presentó una curvatura convexa creciente. El módulo de Young promedio (módulo secante) para la fase de carga fue $Y_l = \qty{538 \pm 25}{kPa}$, valor representativo para cauchos naturales bajo deformaciones moderadas. Se concluye que el uso del esfuerzo real es indispensable para observar el endurecimiento del material por reducción de área.

    \item \textbf{Cuantificación de la histéresis elástica:}
    Se confirmó la presencia de histéresis mediante la formación de un lazo cerrado en el diagrama de esfuerzo real. La energía mecánica disipada internamente por unidad de volumen durante el ciclo de carga--descarga se estimó en $A = \qty{13{,}817}{kJ/m^3}$, consecuencia de la fricción molecular y la naturaleza viscoelástica de la liga.

    \item \textbf{Estimación del trabajo elástico:}
    El trabajo realizado sobre el resorte para alcanzar la tercera
    configuración de carga, calculado mediante integración numérica (regla trapezoidal), fue $W = \qty{0{,}715}{J}$. A diferencia de la liga, esta energía se considera conservativa dada la ausencia de un lazo de histéresis apreciable en el metal.

\end{enumerate}